\chapter{OOP: Classes, Inheritance, and Method Resolution}
\label{ch:oop}

\section{Class Compilation}

Python classes compile to $\mathsf{Ref}$ objects with shared method
symbols:

\begin{lstlisting}[caption={Class compilation}]
# class MyClass:
#     def __init__(self, x):
#         self.x = x
#     def compute(self):
#         return self.x * 2

# Compilation:
# MyClass(v) => Ref(fresh_addr)
# obj.x      => py_attr_x(obj)
# obj.compute() => py_meth_compute(obj)
#               => inlined: py_attr_x(obj) * 2

def compile_class(cls_node, env, T):
    env.put(cls_node.name, ('__class__', cls_node))
    # Methods are stored for inlining
    for item in cls_node.body:
        if isinstance(item, ast.FunctionDef):
            env.put(f'{cls_node.name}.{item.name}',
                    ('__func__', item))
\end{lstlisting}

\section{Classes as OTerm Inhabitants}

In the OTerm framework, a class instance is not a primitive; it is
an \emph{addressed record}---a reference paired with a mapping from
attribute names to OTerm values.

\begin{definition}[Class OTerm]
Let $C$ be a class with attributes $a_1, \ldots, a_k$ and methods
$m_1, \ldots, m_j$.  An instance $o$ of $C$ compiles to:
\[
  \mathsf{OTerm}_C(o) \coloneqq \bigl(\mathsf{Ref}(\alpha),\;
    \{a_i \mapsto \mathsf{py\_attr}_{a_i}(\alpha)\}_{i=1}^k,\;
    \{m_i \mapsto \lambda \mathsf{self}.\, \mathsf{body}_{m_i}[\mathsf{self}/\alpha]\}_{i=1}^j
  \bigr)
\]
where $\alpha$ is a fresh address constant and each method body is
inlined with $\mathsf{self}$ bound to $\alpha$.
\end{definition}

The key insight is that method calls reduce to ordinary function
application on shared symbols:
$\texttt{obj.compute()} \;\leadsto\; \mathsf{py\_meth\_compute}(\alpha)
  \;\leadsto\; \mathsf{py\_attr\_x}(\alpha) \cdot 2$.
Two programs calling the same method on objects of the same class
produce identical shared-symbol expressions, making equivalence
checking straightforward.

\section{Inheritance as Fiber Refinement}

A class hierarchy $A \leftarrow B \leftarrow C$ (C inherits B
inherits A) defines a \emph{refinement} of the type fibration:
the fiber $\mathsf{TRef}$ is refined into sub-fibers for each class.

\begin{definition}[Class Fiber]
For class $C$, the class fiber $U_C = \{v : \PyObj \mid
\mathsf{tag}(v) = \mathsf{TRef} \land \mathsf{class\_of}(v) = C\}$.

The method resolution order (MRO) determines which method is called:
$\mathsf{meth\_m}(v) = \mathsf{mro}(C)[0].\mathsf{m}(v)$.
\end{definition}

Both programs using the same class hierarchy get the same MRO,
hence the same method dispatch, hence the same shared symbols for
each method call.

\section{Method Resolution as a Path in the Class HIT}

Python's C3 linearisation of the MRO can be modelled as a
\emph{path} in a higher inductive type encoding the class hierarchy.

\begin{definition}[Class Hierarchy HIT]
Given classes $C_1, \ldots, C_n$ with inheritance edges
$C_i \to C_j$ (``$C_i$ inherits from $C_j$''), define the HIT
$\mathsf{ClassHier}$ with:
\begin{itemize}
  \item \textbf{Point constructors.}
    $\mathsf{cls}_i : \mathsf{ClassHier}$ for each class $C_i$.
  \item \textbf{Path constructors.}
    $\mathsf{inherits}_{ij} : \mathsf{cls}_i =_{\mathsf{ClassHier}} \mathsf{cls}_j$
    for each inheritance edge $C_i \to C_j$.
\end{itemize}
\end{definition}

A method lookup for method $m$ on an object of class $C_i$ is a
\emph{transport} along the MRO path:
\[
  \mathsf{resolve}_m(C_i) \coloneqq
  \mathsf{transp}^{\mathsf{HasMethod}_m}(\mathsf{mro\_path}(C_i))
\]
where $\mathsf{HasMethod}_m$ is the type family assigning to each
class its implementation of $m$ (or $\bot$ if absent), and
$\mathsf{mro\_path}(C_i)$ is the composite path through the
linearised ancestor chain.

\begin{proposition}[MRO Determinism]
For a fixed class hierarchy, the MRO path is unique (up to
homotopy), so $\mathsf{resolve}_m(C_i)$ is well-defined.
Two programs with the same class hierarchy resolve every method
identically.
\end{proposition}

\section{Virtual Dispatch as Fiber Case Split}

When a function receives an argument of an abstract base type and
calls a virtual method, the compiler performs a \emph{case split on
the object's class fiber}:

\begin{lstlisting}[caption={Virtual dispatch compilation}]
# obj.area() where obj : Shape (base class)
# compile to:
If(class_of(obj) == Circle,
   py_meth_area_Circle(obj),
   If(class_of(obj) == Rect,
      py_meth_area_Rect(obj),
      py_meth_area_Shape(obj)))   # default
\end{lstlisting}

In sheaf-theoretic terms, the class tag induces a covering
$\{U_{C_1}, \ldots, U_{C_n}\}$ of the $\mathsf{TRef}$ fiber,
and the virtual dispatch is exactly the local section assignment
on each sub-cover.  If two programs dispatch the same set of
methods on the same class hierarchy, their local sections agree
on each $U_{C_i}$, producing the same global section and hence
$\Hone = 0$.

\section{Descriptors and Properties}

$\texttt{@property}$ compiles to a shared function:
$\mathsf{obj.prop}$ $\to$ $\mathsf{py\_prop\_prop}(\mathsf{obj})$.
Since both programs access the same property, they get the same
shared symbol.

More generally, Python's descriptor protocol (\texttt{\_\_get\_\_},
\texttt{\_\_set\_\_}, \texttt{\_\_delete\_\_}) compiles to shared
functions parameterised by the descriptor's owner class and the
accessing instance.  Two programs using the same descriptor on the
same class produce identical OTerm expressions.


\chapter{Error Handling and Graceful Degradation}
\label{ch:errors}

\section{Degradation Cascade}

The pipeline degrades through three levels of decreasing certainty:

\begin{center}
\begin{tabular}{c l l}
  \textbf{Level} & \textbf{Trigger} & \textbf{Verdict} \\
  \hline
  1 & Compiler succeeds, Z3 decides &
      $\mathsf{EQUIVALENT}$ or $\mathsf{INEQUIVALENT}$ \\
  2 & Compiler succeeds, Z3 times out &
      $\mathsf{UNKNOWN}$ (partial confidence) \\
  3 & Compiler fails on unsupported syntax &
      $\mathsf{INCONCLUSIVE}$ (fresh-constant fallback) \\
\end{tabular}
\end{center}

\noindent
Level~1 is the normal case: the AST compiler produces OTerms, Z3
decides per-fiber equivalence, and the \v{C}ech complex yields a
definitive verdict.  Level~2 retains soundness---no false
equivalences are reported---but loses completeness on timed-out
fibers.  Level~3 is the most conservative: the fresh-constant
strategy guarantees soundness at the cost of systematic false
negatives.

\begin{proposition}[Soundness of Degradation]
At every level, if the pipeline reports $\mathsf{EQUIVALENT}$,
then the two programs are semantically equivalent on all verified
fibers.  No degradation level can produce a false positive.
\end{proposition}

\section{Unsupported Syntax}

When the AST compiler encounters a construct it cannot handle
(e.g., \texttt{async/await}, \texttt{match/case}, \texttt{global},
\texttt{exec}), it produces a \emph{fresh uninterpreted constant}:
$\mathsf{FreshConst}(\PyObj, \texttt{unsupported})$.

This is \emph{sound}: the fresh constant is unique, so two programs
with different unsupported constructs produce different terms (no
false equivalences).  It is \emph{incomplete}: two programs with the
\emph{same} unsupported construct also produce different fresh
constants (false inequivalence).

\section{Z3 Timeout Handling}

When Z3 times out on a fiber:
\begin{enumerate}
  \item Mark the fiber as $\mathsf{UNKNOWN}$.
  \item Continue checking other fibers.
  \item In the $\Hone$ computation, treat UNKNOWN fibers as
    ``not verified'' (not as failures).
  \item Report a partial verdict: ``Equivalent on $k/n$ fibers,
    $m$ fibers timed out.''
\end{enumerate}

\begin{lstlisting}[caption={Timeout handling with escalation}]
def check_with_timeout(fiber, secs_f, secs_g, solver, timeout):
    solver.set('timeout', timeout)
    solver.push()
    # ... add constraints ...
    result = solver.check()
    solver.pop()
    if result == unknown:
        # Escalate: retry with 3x timeout on just this fiber
        solver.set('timeout', timeout * 3)
        result = solver.check()
    return result
\end{lstlisting}

\section{Per-Fiber Timeouts and Global Deadlines}

The benchmark harness enforces a two-level timeout architecture:

\begin{enumerate}
  \item \textbf{Per-fiber Z3 timeout} (default 5\,000\,ms).
    Set via \texttt{z3.set\_param('timeout', 5000)}.  Controls
    the maximum time Z3 spends on any single satisfiability query.
    Fibers exceeding this limit are marked $\mathsf{UNKNOWN}$ and
    the pipeline continues.

  \item \textbf{Global subprocess deadline} (default 12\,s wall-clock).
    Each benchmark pair runs in a separate subprocess with a hard
    \texttt{subprocess.run(timeout=12)} deadline.  If the entire
    pipeline (compilation $+$ all fiber checks $+$ $\Hone$) exceeds
    this budget, the subprocess is killed and the pair is marked
    $\mathsf{TIMEOUT}$.
\end{enumerate}

\noindent
The per-fiber timeout provides fine-grained control: if fiber $U_3$
times out but fibers $U_1, U_2, U_4$ succeed, the pipeline still
reports a partial verdict.  The global deadline is a safety net
against pathological cases (e.g., exponentially many fibers from
duck type inference).

\section{Memory Safety}

Each subprocess is sandboxed with OS-level resource limits to
prevent a single hard pair from crashing the host:

\begin{lstlisting}[caption={Resource limits from \texttt{run\_benchmark.py}}]
import resource
# Virtual memory: 768 MB hard limit
resource.setrlimit(resource.RLIMIT_AS,
                   (768*1024*1024, 768*1024*1024))
# CPU time: 10s soft / 12s hard (SIGXCPU / SIGKILL)
resource.setrlimit(resource.RLIMIT_CPU, (10, 12))
# Z3 internal memory cap
z3.set_param('memory_max_size', 512)   # 512 MB
\end{lstlisting}

\noindent
The three limits interact as a defence-in-depth stack:
\begin{itemize}
  \item \textbf{Z3 memory\_max\_size} (512\,MB) is the innermost
    guard.  If Z3's internal allocator exceeds this, the solver
    returns $\mathsf{unknown}$ rather than allocating more memory.
  \item \textbf{RLIMIT\_AS} (768\,MB) catches memory growth outside
    Z3 (e.g., large AST compilation, Python object overhead).
    Exceeding it raises \texttt{MemoryError}.
  \item \textbf{RLIMIT\_CPU} (10/12\,s) catches infinite loops in
    compilation or in Z3's E-matching engine.  The soft limit
    sends \texttt{SIGXCPU}; the hard limit sends \texttt{SIGKILL}.
\end{itemize}

\section{Combinatorial Explosion in Site Category}

When duck type inference allows $k$ types per parameter and there
are $n$ parameters, the number of fiber combinations is $k^n$.
Mitigation:
\begin{enumerate}
  \item Cap at 50 fiber combinations (skip rare combinations).
  \item Prioritize fibers by duck type confidence (high-confidence
    single-type fibers first).
  \item Use the spectral sequence (Chapter~\ref{ch:spectral}) to
    check important fibers first and skip if degeneration is detected.
\end{enumerate}

\section{Partial Verdicts and Confidence Levels}

\begin{definition}[Confidence Level]
$\mathsf{confidence}(f, g) = \frac{\text{fibers verified equivalent}}
{\text{total non-timed-out fibers}}$

A confidence of 1.0 means full equivalence proof.
A confidence of 0.8 means 80\% of fibers verified, 20\% timed out.
A confidence of 0.0 with counterexample means definite inequivalence.
\end{definition}


\chapter{Scalability}
\label{ch:scalability}

\section{Complexity of the Full Pipeline}

\begin{center}
\begin{tabular}{l l l}
  \textbf{Stage} & \textbf{Complexity} & \textbf{Bottleneck} \\
  \hline
  AST parsing & $O(|S|)$ & Python \texttt{ast.parse} \\
  AST compilation & $O(|A|)$ & single-pass OTerm build \\
  Duck type inference & $O(|A|)$ & single AST walk \\
  Site construction & $O(k^n)$ & Cartesian product of types \\
  Per-fiber Z3 check & $O(k^n \cdot C_\mathsf{Z3})$ &
    \textbf{dominant term} \\
  $\Hone$ computation & $O(k^{3n})$ & GF(2) matrix rank \\
  \hline
  \textbf{Total} & $O(k^n \cdot C_\mathsf{Z3} + k^{3n})$ & \\
\end{tabular}
\end{center}

For the common case ($k = 1$, single-typed parameters):
$O(C_\mathsf{Z3})$ --- a single Z3 query.

For the benchmark ($k \leq 2$, $n \leq 4$):
$O(16 \cdot C_\mathsf{Z3} + 4096)$ --- at most 16 Z3 queries plus
cheap linear algebra.  Runs in $< 3$ seconds.

\section{Bottleneck Analysis}

The dominant cost is the per-fiber Z3 check:
$O(k^n \cdot C_\mathsf{Z3})$.  This product has two axes of growth:

\begin{enumerate}
  \item \textbf{Number of fibers} ($k^n$).
    Controlled by duck type inference precision.  In practice,
    most parameters have $k = 1$ (monomorphic); polymorphic
    parameters with $k = 2$ and $n \leq 4$ yield at most 16
    fibers.

  \item \textbf{Per-fiber solver time} ($C_\mathsf{Z3}$).
    Depends on term depth, axiom count, and use of recursive
    functions (see below).  Ground queries finish in 1--50\,ms;
    queries with RecFunctions take 50--500\,ms; queries with
    quantified axioms may reach 500\,ms--5\,s.
\end{enumerate}

\noindent
The $\Hone$ computation ($O(k^{3n})$, Gaussian elimination over
GF(2)) is never the bottleneck: for $k^n = 16$ fibers, the
coboundary matrix is $16 \times 16$ and rank computation takes
microseconds.

\section{Z3 Query Complexity}

The Z3 query complexity $C_\mathsf{Z3}$ depends on:
\begin{itemize}
  \item \textbf{Term depth}: deeper Z3 terms (from deeper nesting or
    longer loops) produce harder queries.  Fold RecFunctions add
    depth proportional to the fold range.
  \item \textbf{Number of sections}: more branches produce more
    section pairs to compare.  $m$ sections in A $\times$ $n$
    sections in B $= m \cdot n$ comparisons.
  \item \textbf{Axiom count}: more ForAll axioms slow down E-matching.
    We add axioms only for shared functions actually used (axiom
    relevance filtering).
\end{itemize}

Typical Z3 query time: 1--50ms for ground queries (no RecFunctions),
50--500ms with RecFunctions, 500ms--5s with axioms.

\section{Scaling Strategies}

\begin{enumerate}
  \item \textbf{Fiber pruning.}
    Duck type inference assigns a confidence $p_i \in [0, 1]$ to each
    type assignment.  Fibers whose joint confidence
    $\prod_j p_{i_j} < \epsilon$ (default $\epsilon = 0.01$)
    are pruned before Z3 is invoked, reducing the effective fiber
    count from $k^n$ to a small subset.

  \item \textbf{Early termination.}
    For equivalence: if all fibers checked so far are equivalent and
    the spectral sequence (Chapter~\ref{ch:spectral}) has degenerated
    at $E_1$, skip remaining fibers.
    For inequivalence: if \emph{any} fiber produces a counterexample,
    return $\mathsf{INEQUIVALENT}$ immediately.

  \item \textbf{Parallel solving.}
    Z3 queries on distinct fibers are independent and can be
    dispatched to a thread pool.  With $p$ cores, the wall-clock
    time drops to $O(\lceil k^n / p \rceil \cdot C_\mathsf{Z3})$.
    The benchmark harness uses per-subprocess parallelism (one
    pair per process) rather than intra-pair parallelism, but the
    architecture supports both.

  \item \textbf{Incremental checking.}
    When comparing two versions of the same codebase, only functions
    with changed ASTs need re-verification.  Unchanged function
    results are cached and reused
    (Theorem~\ref{thm:compositional}).

  \item \textbf{Timeout budget.}
    A global timeout budget $T$ (default 10\,s) is divided across
    fibers.  Each fiber receives $T / |\mathrm{fibers}|$ initially;
    fibers that finish quickly donate their remaining budget to
    harder fibers.
\end{enumerate}

\section{Heuristics for Large Programs}

For programs with 100+ lines:
\begin{enumerate}
  \item \textbf{Function-level decomposition}: verify each function
    independently, then compose results (Theorem~\ref{thm:compositional}).
  \item \textbf{Incremental checking}: check modified functions only,
    reuse cached results for unchanged ones.
  \item \textbf{Early termination}: if any fiber produces a
    counterexample, return INEQUIVALENT immediately (no need to
    check remaining fibers).
  \item \textbf{Timeout budget}: allocate a total timeout budget
    (e.g., 10s) and divide across fibers.  Skip low-priority fibers
    if budget is exhausted.
\end{enumerate}


% ══════════════════════════════════════════════════════════
% APPENDICES
% ══════════════════════════════════════════════════════════

\appendix

\chapter{Mathematical Prerequisites}
\label{app:prereqs}

This appendix collects the mathematical background used throughout
the monograph.  We state definitions and key theorems without proof;
references to standard texts are given at the end of each section.

\section{Category Theory Essentials}

A \emph{category} $\mathcal{C}$ consists of objects, morphisms,
composition, and identity, satisfying associativity and unit laws.
A \emph{functor} $F : \mathcal{C} \to \mathcal{D}$ preserves
composition and identity.  A \emph{natural transformation}
$\eta : F \Rightarrow G$ is a family of morphisms $\eta_A : F(A) \to
G(A)$ commuting with all morphisms.

Key concepts used:
\begin{itemize}
  \item \textbf{Initial algebra} (Ch.~\ref{ch:nat-elim},
    \ref{ch:recursion-zoo}).
    The initial object in the category of $F$-algebras for an
    endofunctor $F$.  Lists are the initial algebra for
    $F(X) = 1 + A \times X$.  The unique morphism from the
    initial algebra is the \emph{catamorphism} (fold).

  \item \textbf{Grothendieck construction}
    (Ch.~\ref{ch:sheaves}, \ref{ch:duck-detail}).
    Turns a pseudofunctor
    $F : \mathcal{C}^{\mathrm{op}} \to \Cat$ into a fibered category
    $\int F$ over $\mathcal{C}$.  Our type fibration is a Grothendieck
    construction: objects of $\int F$ are pairs $(U, s)$ where $U$ is
    a fiber and $s \in F(U)$ is a local section.

  \item \textbf{Yoneda lemma} (Ch.~\ref{ch:enrichment}).
    $\mathrm{Nat}(\mathrm{Hom}(A, -), F)
    \cong F(A)$.  Functions on representable presheaves are determined
    by their value at the representing object.

  \item \textbf{Adjunctions.}
    A pair of functors $F \dashv G$ with natural bijection
    $\mathrm{Hom}(FA, B) \cong \mathrm{Hom}(A, GB)$.
    Used implicitly in the compilation (syntax $\dashv$ semantics).

  \item \textbf{Limits and colimits.}
    Products, pullbacks, pushouts, and coequalisers appear in the
    site category construction and the \v{C}ech complex.
\end{itemize}

\noindent
\textbf{References.}
Mac~Lane, \emph{Categories for the Working Mathematician}
\cite{maclane1971};
Riehl, \emph{Category Theory in Context} \cite{riehl2017}.

\section{Sheaf Theory and Cohomology}

A \emph{site} is a category with a Grothendieck topology (covering
families).  A \emph{presheaf} is a contravariant functor to $\Set$.
A \emph{sheaf} satisfies the gluing axiom: local sections that agree
on overlaps glue uniquely.

The \emph{\v{C}ech complex} computes cohomology from a covering
$\mathcal{U} = \{U_i\}$:
\[
  C^p(\mathcal{U}, \presheaf) \coloneqq
  \prod_{i_0 < \cdots < i_p} \presheaf(U_{i_0} \cap \cdots \cap U_{i_p})
\]
with alternating coboundary operators $\delta^p : C^p \to C^{p+1}$.

\begin{itemize}
  \item $\Hzero(\mathcal{U}, \presheaf) = \ker \delta^0$ is the space of
    \emph{global sections}---type assignments that are consistent
    across all fibers.
  \item $\Hone(\mathcal{U}, \presheaf) = \ker \delta^1 / \mathrm{im}\, \delta^0$
    measures \emph{obstruction to gluing}: a non-trivial $\Hone$ class
    means local sections cannot be reconciled into a global one.
    In our framework, $\Hone \neq 0$ signals a semantic bug.
\end{itemize}

\noindent
\textbf{References.}
Grothendieck, \emph{Sur quelques points d'alg\`ebre homologique}
\cite{grothendieck1957};
Hartshorne, \emph{Algebraic Geometry}, Ch.~III \cite{hartshorne1977};
Curry, \emph{Sheaves, Cosheaves and Applications} \cite{curry2014}.

\section{Homotopy Type Theory Essentials}

Types are spaces.  Terms are points.  Equalities are paths.
\[
  \Path_A(a, b) = \{p : \interval \to A \mid p(0) = a,\; p(1) = b\}
\]
Function extensionality: $(f = g) \simeq \prod_x f(x) = g(x)$.

\noindent
Univalence: $(A \simeq B) \simeq (A =_\Type B)$.  Equivalent types
are identical, which in our framework justifies treating two programs
with isomorphic OTerm representations as equal.

\noindent
Higher inductive types (HITs): types with point constructors \emph{and}
path constructors.  Used for:
\begin{itemize}
  \item The interval type $\interval$ (Chapter~\ref{ch:interval}).
  \item Pushouts and coequalisers in the site category
    (Chapter~\ref{ch:cech}).
  \item The class hierarchy HIT (Chapter~\ref{ch:oop}).
\end{itemize}

\noindent
\textbf{References.}
The Univalent Foundations Program, \emph{Homotopy Type Theory} \cite{hott};
Cohen et al., \emph{Cubical Type Theory} \cite{cohen2018};
Angiuli et al., \emph{Computational Higher-Dimensional Type Theory}
\cite{angiuli2018}.

\section{Algebraic Geometry Background}

The algebraic-geometric concepts used in this monograph (beyond
sheaf cohomology) include:

\begin{itemize}
  \item \textbf{Fibrations and fibers.}
    The type fibration $\pi : \int \mathsf{Ty} \to \site$ assigns
    to each open set $U$ in the site a fiber of possible type
    assignments.  Local sections are type-annotated program fragments
    over $U$.

  \item \textbf{Spectral sequences} (Chapter~\ref{ch:spectral}).
    The Mayer--Vietoris spectral sequence
    $E_1^{p,q} \Rightarrow H^{p+q}$ provides a filtration of the
    cohomology that allows early termination: if $E_1$ degenerates
    (all differentials vanish), then $\Hone = 0$ without computing
    the full \v{C}ech complex.

  \item \textbf{Descent data.}
    A collection of local sections with transition isomorphisms
    on overlaps.  The sheaf condition says that descent data is
    \emph{effective}---it glues to a unique global section.
    Failure of effectiveness is precisely $\Hone \neq 0$.
\end{itemize}

\noindent
\textbf{References.}
Hartshorne, \emph{Algebraic Geometry} \cite{hartshorne1977};
Grothendieck, \emph{SGA 4} (for sites and topoi).

\section{SMT Solving Essentials}

Z3 decides satisfiability of first-order formulas in combined theories
(arithmetic, strings, arrays, algebraic datatypes).  Key features
used:
\begin{itemize}
  \item \textbf{Algebraic datatypes}: the $\PyObj$ sort is a Z3
    datatype with constructors for each Python type tag
    (Chapter~\ref{ch:pyobj}).
  \item \textbf{Recursive functions}: fold/map/filter compile to Z3
    \texttt{RecFunction} declarations, enabling bounded reasoning
    about iteration (Chapter~\ref{ch:recursion-zoo}).
  \item \textbf{Quantifiers}: ForAll axioms express shared-function
    semantics.  Z3 uses E-matching and MBQI for instantiation
    (Chapter~\ref{ch:axioms}).
  \item \textbf{Enumeration sorts}: used for type tags in the duck
    type lattice (Chapter~\ref{ch:duck-detail}).
  \item \textbf{Timeouts}: per-query timeouts and memory limits
    ensure termination (Chapter~\ref{ch:errors}).
\end{itemize}

\noindent
\textbf{References.}
de~Moura and Bj{\o}rner, \emph{Z3: An Efficient SMT Solver}
\cite{moura2008};
Barrett et al., \emph{The SMT-LIB Standard} (for theory definitions).


\chapter{Benchmark Pair Catalog}
\label{app:pairs}

The complete catalog of 50 equivalent and 50 non-equivalent benchmark
pairs is maintained in the repository at
\texttt{benchmarks/hard\_equiv\_suite.py} and
\texttt{tests/test\_equivalence/test\_hard\_eq\_pairs.py}.

Each pair consists of two Python functions (20--60 lines each) that
exercise deep Python semantics.  The equivalent pairs span all 24
dichotomies; the non-equivalent pairs test subtle behavioural
differences including mutation, aliasing, closure binding, NaN
handling, and edge cases.

\section{Pipeline Resolution Stages}

Each pair is resolved by a specific stage of the pipeline.  We
classify the resolution stage for each pair:

\begin{center}
\begin{tabular}{l p{8cm}}
  \textbf{Stage} & \textbf{Description} \\
  \hline
  \textsf{Compile} &
    Both programs compile to identical or $\alpha$-equivalent OTerms.
    Z3 confirms $\forall x.\, f(x) = g(x)$ in a single ground query.
    \\
  \textsf{Fiber} &
    Programs differ syntactically but agree on each fiber after
    per-fiber Z3 checking.  Resolved by the local check stage. \\
  \textsf{Cohomology} &
    Local sections disagree on some fiber overlaps; the
    $\Hone$ computation determines whether disagreements are
    coboundaries (equivalent) or cocycles (inequivalent). \\
  \textsf{Axiom} &
    Equivalence depends on shared-function axioms (e.g., commutativity
    of addition, associativity of string concatenation).  Resolved
    by Z3 with ForAll axioms. \\
  \textsf{RecFunction} &
    Programs use different iteration patterns (fold vs.\ recursion
    vs.\ comprehension).  Resolved by Z3 RecFunction reasoning. \\
\end{tabular}
\end{center}

\section{Difficulty Classification}

We classify all 100 pairs into three difficulty tiers:

\begin{description}
  \item[Trivial] (canonical match).
    Both programs use the same algorithm with minor syntactic
    variation (variable renaming, loop $\leftrightarrow$
    comprehension, \texttt{sorted()} $\leftrightarrow$ manual sort).
    Resolved at the Compile or single-fiber stage in $< 50$\,ms.

  \item[Medium] (path search).
    Programs use different data structures or iteration orders but
    compute the same function.  Requires multi-fiber checking and/or
    axiom-assisted reasoning.  Typical resolution time: 50--500\,ms.

  \item[Hard] (different algorithms).
    Programs implement fundamentally different algorithms for the
    same problem (e.g., recursive descent vs.\ shunting-yard parser,
    convex hull via gift-wrapping vs.\ Graham scan, BFS vs.\ DFS
    with post-order).  Requires RecFunction reasoning, multiple
    axioms, or cohomological analysis.  Resolution time: 500\,ms--5\,s.
\end{description}

\section{Equivalent Pairs (EQ-01 through EQ-50)}

\begin{center}
\small
\begin{tabular}{r p{5.5cm} l l}
  \textbf{\#} & \textbf{Description} & \textbf{Stage} & \textbf{Tier} \\
  \hline
  01 & Flatten nested data: recursive generator vs.\ iterative stack
     & RecFunction & Hard \\
  02 & Word frequency: regex+Counter vs.\ char state machine+heap
     & RecFunction & Hard \\
  03 & Expression evaluator: recursive vs.\ stack-based
     & RecFunction & Hard \\
  04 & Register VM: dict-based vs.\ list-based
     & Axiom & Medium \\
  05 & Topological sort: Kahn's vs.\ DFS post-order
     & RecFunction & Hard \\
  06 & Run-length encoding: groupby vs.\ manual loop
     & Fiber & Medium \\
  07 & Convex hull: gift-wrapping vs.\ Graham scan
     & RecFunction & Hard \\
  08 & Edit distance: full DP matrix vs.\ rolling two-row
     & RecFunction & Hard \\
  09 & LRU cache: OrderedDict vs.\ dict+doubly-linked list
     & Axiom & Hard \\
  10 & Matrix multiply: nested loops vs.\ zip+map
     & Fiber & Medium \\
  11 & Merge sorted iterables: heapq vs.\ manual merge
     & RecFunction & Hard \\
  12 & Binary tree traversal: recursive vs.\ iterative stack
     & RecFunction & Medium \\
  13 & Max flow: Ford--Fulkerson BFS vs.\ DFS path
     & RecFunction & Hard \\
  14 & Arithmetic parser: recursive descent vs.\ shunting-yard
     & RecFunction & Hard \\
  15 & Quickselect vs.\ sort-and-index
     & RecFunction & Hard \\
  16 & Strongly connected components: Tarjan vs.\ Kosaraju
     & RecFunction & Hard \\
  17 & Power set: recursive vs.\ bitmask enumeration
     & RecFunction & Medium \\
  18 & Pretty-print nested data: recursive vs.\ stack+indent
     & RecFunction & Medium \\
  19 & 0/1 knapsack: DP table vs.\ memoised recursion
     & RecFunction & Hard \\
  20 & Deep copy: recursive vs.\ \texttt{copy.deepcopy}
     & Axiom & Medium \\
  21 & Prime factorisation: trial division vs.\ sieve
     & RecFunction & Medium \\
  22 & Cycle detection: Floyd's vs.\ visited-set DFS
     & RecFunction & Hard \\
  23 & Bracket matching: stack vs.\ counter with type check
     & Fiber & Trivial \\
  24 & Levenshtein distance: full matrix vs.\ two-row DP
     & RecFunction & Hard \\
  25 & Interval merge: sort+scan vs.\ priority queue
     & Fiber & Medium \\
  26 & $k$-th smallest: quickselect vs.\ heapq.nsmallest
     & RecFunction & Medium \\
  27 & LCS: DP table vs.\ memoised recursion
     & RecFunction & Hard \\
  28 & Regex match: NFA simulation vs.\ recursive backtrack
     & RecFunction & Hard \\
  29 & Graph inversion: comprehension vs.\ defaultdict loop
     & Fiber & Trivial \\
  30 & Prefix search: trie vs.\ sorted+bisect
     & RecFunction & Hard \\
  31 & Dijkstra: heapq vs.\ manual priority queue
     & RecFunction & Hard \\
  32 & Island counting: BFS vs.\ DFS flood fill
     & RecFunction & Hard \\
  33 & Permutation by index: factorial number system vs.\ itertools
     & Axiom & Medium \\
  34 & RPN evaluator: stack vs.\ recursive
     & Fiber & Medium \\
  35 & Take-while: itertools vs.\ manual loop
     & Fiber & Trivial \\
  36 & Dict inversion: comprehension vs.\ setdefault loop
     & Fiber & Trivial \\
  37 & Fibonacci: matrix exponentiation vs.\ iterative
     & RecFunction & Hard \\
  38 & Anagram grouping: sorted-key vs.\ Counter-key
     & Axiom & Medium \\
  39 & Perfect power check: log-based vs.\ brute-force
     & RecFunction & Medium \\
  40 & Binomial coefficient: multiplicative vs.\ Pascal's triangle
     & RecFunction & Medium \\
  41 & Zip longest: manual fill vs.\ itertools
     & Fiber & Trivial \\
  42 & Memoised function application: decorator vs.\ manual cache
     & Axiom & Medium \\
  43 & CSV parse: split-based vs.\ state machine
     & RecFunction & Medium \\
  44 & Coin change: DP vs.\ BFS
     & RecFunction & Hard \\
  45 & GCD: Euclidean vs.\ subtraction
     & RecFunction & Medium \\
  46 & Lisp interpreter: recursive eval vs.\ CPS
     & RecFunction & Hard \\
  47 & Stock trading: DP vs.\ state machine
     & RecFunction & Hard \\
  48 & Partition function: generating function vs.\ DP
     & RecFunction & Hard \\
  49 & Deep equality: recursive vs.\ stack-based
     & RecFunction & Medium \\
  50 & Catalan number: closed-form vs.\ DP
     & RecFunction & Medium \\
\end{tabular}
\end{center}

\section{Non-Equivalent Pairs (NEQ-01 through NEQ-50)}

\begin{center}
\small
\begin{tabular}{r p{5cm} l l}
  \textbf{\#} & \textbf{Witness Divergence} & \textbf{Stage} & \textbf{Tier} \\
  \hline
  01 & Floor div vs.\ truncation on negative &
       Compile & Trivial \\
  02 & Python \texttt{\%} vs.\ C-style \texttt{\%} &
       Compile & Trivial \\
  03 & Banker's rounding vs.\ round-half-up &
       Fiber & Medium \\
  04 & \texttt{split()} vs.\ \texttt{split(' ')} &
       Axiom & Medium \\
  05 & \texttt{==} vs.\ \texttt{is} for large integers &
       Compile & Trivial \\
  06 & Late-binding vs.\ early-binding closure &
       RecFunction & Hard \\
  07 & Mutable default argument accumulation &
       Axiom & Hard \\
  08 & Insertion order vs.\ sorted dict keys &
       Fiber & Medium \\
  09 & NaN propagation in \texttt{max()} &
       Fiber & Medium \\
  10 & \texttt{or} returns value vs.\ \texttt{bool} &
       Compile & Trivial \\
  11 & Float accumulation order &
       Fiber & Hard \\
  12 & Sorted-set dedup vs.\ insertion-order dedup &
       Fiber & Medium \\
  13 & \texttt{except Exception} vs.\ bare \texttt{except} &
       Axiom & Hard \\
  14 & Stable sort vs.\ key-only sort &
       Fiber & Medium \\
  15 & List vs.\ generator double-iteration &
       RecFunction & Hard \\
  16 & Floor div vs.\ true div result type &
       Compile & Trivial \\
  17 & List \texttt{count} vs.\ set membership &
       Fiber & Medium \\
  18 & \texttt{str()} vs.\ \texttt{repr()} &
       Axiom & Medium \\
  19 & Dict merge order with conflicts &
       Fiber & Medium \\
  20 & Truthiness vs.\ \texttt{None} check for empty list &
       Compile & Trivial \\
  21 & Chained vs.\ parenthesised comparison &
       Compile & Trivial \\
  22 & \texttt{+=} vs.\ \texttt{+} on list aliases &
       Axiom & Hard \\
  23 & Float \texttt{sqrt} vs.\ \texttt{isqrt} precision &
       Fiber & Hard \\
  24 & List comp vs.\ generator exception location &
       RecFunction & Hard \\
  25 & Stable ascending vs.\ reverse-reverse sort &
       Fiber & Medium \\
  26 & \texttt{any} short-circuit vs.\ full-scan count &
       Fiber & Medium \\
  27 & Walrus operator scope &
       Compile & Medium \\
  28 & \texttt{dict.get} vs.\ \texttt{KeyError} handling &
       Axiom & Medium \\
  29 & \texttt{[0]*n} shallow vs.\ comprehension deep copy &
       Axiom & Hard \\
  30 & \texttt{list} vs.\ \texttt{map} type name &
       Compile & Trivial \\
  31 & \texttt{enumerate} from 0 vs.\ from 1 &
       Compile & Trivial \\
  32 & \texttt{zip} truncate vs.\ \texttt{zip\_longest} &
       Fiber & Medium \\
  33 & \texttt{sorted} returns list vs.\ \texttt{sort} returns None &
       Compile & Trivial \\
  34 & String \texttt{a+b} vs.\ \texttt{b+a} &
       Compile & Trivial \\
  35 & Exception in \texttt{try} body vs.\ \texttt{else} clause &
       Axiom & Medium \\
  36 & Counter subtract keeps negatives vs.\ drops &
       Axiom & Medium \\
  37 & Class attribute shared vs.\ instance attribute &
       Axiom & Hard \\
  38 & \texttt{*args} as tuple vs.\ list &
       Compile & Trivial \\
  39 & \texttt{int()} truncation vs.\ \texttt{math.floor} &
       Compile & Trivial \\
  40 & Logical \texttt{and} vs.\ bitwise \texttt{\&} &
       Compile & Trivial \\
  41 & \texttt{range} excludes vs.\ includes endpoint &
       Compile & Trivial \\
  42 & Shallow vs.\ deep copy inner mutation &
       Axiom & Hard \\
  43 & \texttt{==} vs.\ \texttt{is} for empty list &
       Compile & Trivial \\
  44 & \texttt{for/else} with vs.\ without \texttt{break} &
       Fiber & Medium \\
  45 & \texttt{replace} all vs.\ replace first &
       Axiom & Medium \\
  46 & \texttt{dict.update} mutates vs.\ \texttt{\{**d\}} doesn't &
       Axiom & Hard \\
  47 & Parenthesised expression vs.\ tuple &
       Compile & Trivial \\
  48 & \texttt{abs} handles complex; manual doesn't &
       Fiber & Medium \\
  49 & \texttt{max} first vs.\ last on tie &
       Fiber & Medium \\
  50 & Return string vs.\ \texttt{print} returns None &
       Compile & Trivial \\
\end{tabular}
\end{center}

\section{Summary Statistics}

\begin{center}
\begin{tabular}{l r r r}
  & \textbf{Trivial} & \textbf{Medium} & \textbf{Hard} \\
  \hline
  EQ pairs  & 5  & 18 & 27 \\
  NEQ pairs & 19 & 17 & 14 \\
  \hline
  Total     & 24 & 35 & 41 \\
\end{tabular}
\end{center}

\noindent
The distribution is intentionally skewed toward hard pairs: the
benchmark's purpose is to stress-test the pipeline on cases where
naive syntactic comparison fails.  The NEQ pairs are skewed toward
trivial/medium because many Python semantic traps (NaN, mutable
defaults, \texttt{is} vs.\ \texttt{==}) produce divergence that
the compiler detects early.


% ══════════════════════════════════════════════════════════
% BIBLIOGRAPHY
% ══════════════════════════════════════════════════════════

\begin{thebibliography}{99}

\bibitem{cohen2018}
C.~Cohen, T.~Coquand, S.~Huber, A.~M\"ortberg.
\emph{Cubical Type Theory: A Constructive Interpretation of the
  Univalence Axiom}.
Journal of Automated Reasoning, 2018.

\bibitem{cchm2017}
T.~Coquand, S.~Huber, A.~M\"ortberg.
\emph{On Higher Inductive Types in Cubical Type Theory}.
LICS 2018.

\bibitem{hott}
The Univalent Foundations Program.
\emph{Homotopy Type Theory: Univalent Foundations of Mathematics}.
Institute for Advanced Study, 2013.

\bibitem{angiuli2018}
C.~Angiuli, R.~Harper, T.~Wilson.
\emph{Computational Higher-Dimensional Type Theory}.
POPL 2018.

\bibitem{voevodsky2010}
V.~Voevodsky.
\emph{Univalent Foundations of Mathematics}.
WoLLIC 2010.

\bibitem{moura2008}
L.~de~Moura, N.~Bj{\o}rner.
\emph{Z3: An Efficient SMT Solver}.
TACAS 2008.

\bibitem{politz2013}
J.~G.~Politz et al.
\emph{Python: The Full Monty}.
OOPSLA 2013.

\bibitem{grothendieck1957}
A.~Grothendieck.
\emph{Sur quelques points d'alg\`ebre homologique}.
T\^ohoku Math.\ J., 1957.

\bibitem{reynolds1972}
J.~C.~Reynolds.
\emph{Definitional Interpreters for Higher-Order Programming Languages}.
ACM Annual Conference, 1972.

\bibitem{pnueli1998}
A.~Pnueli, M.~Siegel, E.~Singerman.
\emph{Translation Validation}.
TACAS 1998.

\bibitem{godlin2009}
B.~Godlin, O.~Strichman.
\emph{Regression Verification}.
DAC 2009.

\bibitem{barthe2011}
G.~Barthe, J.~M.~Crespo, C.~Kunz.
\emph{Relational Verification Using Product Programs}.
FM 2011.

\bibitem{cousot1977}
P.~Cousot, R.~Cousot.
\emph{Abstract Interpretation}.
POPL 1977.

\bibitem{rondon2008}
P.~M.~Rondon, M.~Kawaguchi, R.~Jhala.
\emph{Liquid Types}.
PLDI 2008.

\bibitem{lumsdaine2011}
P.~L.~Lumsdaine.
\emph{Higher Inductive Types: A Tour}.
2011.

\bibitem{shulman2015}
M.~Shulman.
\emph{Brouwer's Fixed-Point Theorem in Real-Cohesive Homotopy Type Theory}.
MSCS 2018.

\bibitem{curry2014}
J.~Curry.
\emph{Sheaves, Cosheaves and Applications}.
PhD thesis, 2014.

\bibitem{abramsky1987}
S.~Abramsky.
\emph{Domain Theory in Logical Form}.
PhD thesis, 1987.

\bibitem{maclane1971}
S.~Mac~Lane.
\emph{Categories for the Working Mathematician}.
Springer, 1971.

\bibitem{riehl2017}
E.~Riehl.
\emph{Category Theory in Context}.
Dover, 2017.

\bibitem{hartshorne1977}
R.~Hartshorne.
\emph{Algebraic Geometry}.
Springer, 1977.

\end{thebibliography}
